\paragraph{参数敏感性分析}
为检验情景结论对经验参数的依赖程度，对 6 个关键参数分别施加 $\pm 20\%$ 扰动，记录终态江豚 $D$ 和长江鲟 $A$ 相对基线的变化率（表 \ref{tab:q2_sens}）。

\begin{table}[htbp]
\centering
\caption{问题二参数敏感性分析（$\pm 20\%$ 扰动）}
\label{tab:q2_sens}
\begin{tabular}{lcccc}
\toprule
参数 & $\Delta D_{+20\%}$ & $\Delta D_{-20\%}$ & $\Delta A_{+20\%}$ & $\Delta A_{-20\%}$ \\
\midrule
$\delta_B$（通道阻隔） & $-1.9\%$ & $+2.0\%$ & $-2.5\%$ & $+2.6\%$ \\
$R_A$（放流强度） & $-0.2\%$ & $+0.2\%$ & $+10.9\%$ & $-10.9\%$ \\
$e_D$（江豚转化率） & $+8.3\%$ & $-9.4\%$ & $-0.8\%$ & $+0.9\%$ \\
$e_A$（长江鲟转化率） & $-0.0\%$ & $+0.0\%$ & $+1.8\%$ & $-1.8\%$ \\
$a_D$（江豚摄食率） & $-8.5\%$ & $+5.2\%$ & $-2.7\%$ & $+3.6\%$ \\
$a_A$（长江鲟摄食率） & $-0.8\%$ & $+0.8\%$ & $+1.7\%$ & $-1.7\%$ \\
\bottomrule
\end{tabular}
\end{table}

敏感性结果显示，江豚终态对 $e_D$ 和 $a_D$ 最敏感（$|\Delta D|>5\%$），说明江豚恢复对食物获取效率变化最为敏感。长江鲟那边最显眼的是放流强度 $R_A$（$|\Delta A|\approx 11\%$），人工放流仍是恢复中的关键补偿项。通道阻隔 $\delta_B$ 对两物种都有负面影响，只是放到当前参数区间里看，扰动幅度还没有前两者那么大（$<3\%$）。

\begin{figure}[htbp]
\centering
\includegraphics[width=0.88\textwidth]{fig02_q2_rare_species.png}
\caption{问题二和问题五共用的稀有种与入侵模块结果}
\label{fig:q2}
\end{figure}

